\chapter{Blood Clots (Venous Thromboembolism)}

\section*{Definition}
Formation of a blood clot (thrombus) within a deep vein, most commonly in the lower extremities (deep vein thrombosis), which may dislodge and travel to the lungs (pulmonary embolism).

\section*{When to See a Doctor}

\begin{itemize}
 \item Seek urgent medical assessment for one-sided leg swelling, pain, warmth, or redness, which may indicate DVT. Seek emergency care for sudden chest pain, shortness of breath, coughing up blood, fainting, or collapse, which may indicate PE.
\end{itemize}

\section*{How It Develops}
Virchow triad describes predisposing factors: venous stasis (immobility, surgery, long travel), endothelial injury (trauma, surgery, catheterization), and hypercoagulability (inherited thrombophilia, malignancy, pregnancy, oral contraceptives). DVT typically starts in the calf veins and may propagate proximally. Embolization to pulmonary arteries causes ventilation-perfusion mismatch, increased pulmonary vascular resistance, and potentially hemodynamic collapse.

\section*{Causes}
\begin{itemize}
 \item Prolonged immobility (surgery, hospitalization, long-haul travel)
 \item Malignancy (cancer-associated thrombosis)
 \item Oral contraceptives or hormone replacement therapy
 \item Pregnancy and postpartum period
 \item Inherited thrombophilia (Factor V Leiden, prothrombin mutation)
 \item Antiphospholipid syndrome
 \item Trauma or major surgery (especially orthopedic)
 \item Central venous catheters
 \item Obesity
 \item Advanced age
\end{itemize}

\section*{Related Symptoms}
\begin{itemize}
 \item Leg swelling
 \item Chest pain
 \item Shortness of breath
 \item Calf pain
 \item coughing up blood
\end{itemize}


\section*{Medical Evaluation}
\begin{itemize}
 \item A clinician assesses symptoms, vital signs, risk factors, and clinical probability using a validated tool; Homan sign is not reliable for diagnosing DVT.
 \item D-dimer is used selectively in people with low or unlikely clinical probability. A negative result can help exclude VTE in the appropriate pathway, but D-dimer is less specific in pregnancy, cancer, infection, older age, and after surgery.
 \item Compression duplex ultrasound is used for suspected DVT, guided by pre-test probability; CT pulmonary angiography is commonly used for suspected PE when appropriate, with alternative testing when contrast or radiation is unsuitable.
 \item Thrombophilia testing is not routine after every clot. Specialist advice is considered when the result would change management, particularly with a strong family history, unusual thrombosis, or suspected antiphospholipid syndrome.
\end{itemize}

\section*{Treatment Options}
\begin{itemize}
 \item Confirmed DVT or PE is treated with anticoagulation when safe. The choice depends on kidney and liver function, bleeding risk, pregnancy, cancer, drug interactions, and local guidance; direct oral anticoagulants are commonly used in suitable adults.
 \item Low-molecular-weight heparin is commonly used during pregnancy and may be selected in cancer or when oral anticoagulants are unsuitable.
 \item Thrombolysis or catheter-based treatment is reserved for selected life-threatening PE or severe symptomatic iliofemoral DVT under specialist care.
 \item Inferior vena cava filter is placed in patients with contraindications to anticoagulation or recurrent PE despite adequate anticoagulation.
\end{itemize}

\section*{Self-Care and Home Management}
\begin{itemize}
 \item Seek emergency medical care if you have sudden leg swelling, warmth, and redness (possible DVT) or sudden chest pain, shortness of breath, or coughing up blood (possible PE).
 \item Do not massage or apply heat to a swollen, painful leg as this may dislodge a clot and cause pulmonary embolism.
 \item During long journeys, move the legs regularly, walk when practical, and avoid prolonged immobility. People at high risk should ask a clinician about preventive measures rather than self-medicating.
 \item If already diagnosed with a DVT, take anticoagulants exactly as prescribed and follow advice about activity, compression, and follow-up. Do not stop treatment without medical advice.
\end{itemize}

\section*{References}
\begin{thebibliography}{99}
\bibitem{blood_clots:ref1} Di Nisio M, van Es N, et al. ``Deep vein thrombosis and pulmonary embolism.'' \textit{Lancet}. 2016;388(10063):3060-3073.
\bibitem{blood_clots:ref2} Kearon C, Akl EA, et al. ``Antithrombotic therapy for VTE disease: CHEST guideline and expert panel report.'' \textit{Chest}. 2016;149(2):315-352.
\bibitem{blood_clots:ref3} Wolberg AS, Rosendaal FR, et al. ``Venous thrombosis.'' \textit{Nat Rev Dis Primers}. 2015;1:15006.
\bibitem{blood_clots:ref4} Kaushansky K, Lichtman MA, et al. \textit{Williams Hematology}, 10th ed. McGraw-Hill, 2021.
\bibitem{blood_clots:ref5} National Institute for Health and Care Excellence. ``Venous thromboembolic diseases: diagnosis, management and thrombophilia testing (NG158).'' Updated 2023. https://www.nice.org.uk/guidance/ng158.
\end{thebibliography}
